\begin{derivation}[共振单模近场的量子耦合]

一般线性电子--玻色相互作用在相互作用绘景中可写成
\[
\hat V_I(t)
=
\hbar\sum_\lambda
\left[
\hat{\mathcal G}_{\lambda,I}(t)\hat a_\lambda
+
\hat{\mathcal G}_{\lambda,I}^\dagger(t)\hat a_\lambda^\dagger
\right].
\]
这里将光模自由演化的 $\ee^{-\ii\omega_\lambda t}$ 也吸收到 $\hat{\mathcal G}_{\lambda,I}$ 中；$\hat a_\lambda$ 为固定的相互作用绘景模算符，不能再重复附加该相位。$\hat{\mathcal G}_{\lambda,I}$ 仍作用在完整正能电子连续谱上。先把电子自由本征态记为 $|p,s\rangle$，则
\[
\hat{\mathcal G}_{\lambda,I}(t)
=
\sum_{ss'}\int\dd p\,\dd p'\,
\mathcal G_{\lambda;s's}(p',p)\,
\ee^{+\ii\{[E(p')-E(p)]/\hbar-\omega_\lambda\}t}
|p',s'\rangle\langle p,s|.
\]
因此“单模 QPINEM”需要依次做三个投影，而不是一步把完整哈密顿量 改写成升降算符。

第一步是环境谱投影。设目标模式为 $\lambda_0$，其频率为 $\omega_0$；其余模式与目标模式的最小有效失谐记为
\[
\Delta_{\rm spec}
=
\inf_{\lambda\neq\lambda_0}
|\omega_\lambda-\omega_0|.
\]
若作用时间 $T$ 和实际耦合带宽满足
\[
\Delta_{\rm spec}T\gg1,
\qquad
\sup_{\lambda\neq\lambda_0}
\int_{t_i}^{t_f}\dd t\,
\|\hat{\mathcal G}_{\lambda,I}(t)\|
\ll1,
\]
则有限时间 Fourier 核抑制非目标模式，环境和式可约化为 $\lambda_0$。

第二步是电子通道投影。连续谱广义本征态满足 Dirac--$\delta$ 归一化，而离散通道模型需要一组真正 Kronecker 正交的归一态。先在长度 $L_z$ 的周期盒中正则化纵向连续谱，使允许动量间隔为 $2\pi\hbar/L_z$；再选择与主导动量转移 $\hbar\kappa_z$ 相容的一串正交盒归一态。记
\[
|\ell\rangle_V
\equiv
|p_\ell,s_0\rangle_V,
\qquad
p_\ell=p_0+\ell\hbar\kappa_z,
\qquad
{}_V\langle\ell|m\rangle_V=\delta_{\ell m}.
\]
若实际电子是有限宽波包，也可以用彼此近似正交、中心动量为 $p_\ell$ 的归一化通道包代替盒态；最终有效格点只要求相邻通道的重叠误差小于所保留动力学精度。记号上省略盒归一化下标 $V$，而离散化极限保持不变。投影后的电子耦合算符为
\[
\hat{\mathcal G}_{0,I}(t)
=
\sum_\ell
\mathcal G_\ell(t)
\ee^{+\ii\Delta_\ell t}
|\ell+1\rangle\langle\ell|
+
\hat R_{\rm off}(t),
\]
其中
\[
\Delta_\ell
=
\frac{E(p_{\ell+1})-E(p_\ell)}{\hbar}
-\omega_0
\]
是逐边真实失谐，$\hat R_{\rm off}$ 收集自旋翻转、非主导动量转移和其它离散通道。

第三步才是平移对称近似。若显著占据窗口 $|\ell|\le L_{\rm occ}$ 内
\[
\epsilon_{\mathcal G}
\equiv
\max_\ell
\left|
\frac{\mathcal G_\ell-\mathcal G_0}{\mathcal G_0}
\right|
\ll1,
\qquad
\epsilon_\Delta
\equiv
T\max_\ell|\Delta_\ell-\Delta_0|
\ll1,
\]
便可把共同边耦合写成
\[
\mathcal G_\ell(t)\ee^{+\ii\Delta_\ell t}
=
\mathcal G(t)
+
O\!\left(|\mathcal G_0|(\epsilon_{\mathcal G}+\epsilon_\Delta)\right),
\]
并定义双边电子平移算符
\[
\hat b^\dagger
=
\sum_{\ell\in\mathbb Z}
|\ell+1\rangle\langle\ell|,
\qquad
\hat b
=
\sum_{\ell\in\mathbb Z}
|\ell\rangle\langle\ell+1|.
\]
于是
\[
\hat{\mathcal G}_{0,I}(t)
=
\mathcal G(t)\hat b^\dagger
+\delta\hat{\mathcal G}(t).
\]

把这一结果代回单模完整哈密顿量，
\[
\hat V_I(t)
=
\hbar
\left[
\mathcal G(t)\hat b^\dagger\hat a
+
\mathcal G^*(t)\hat b\hat a^\dagger
\right]
+\hat R_{\rm RWA}(t)
+\hat R_{\rm tr}(t),
\]
其中 $\hat R_{\rm tr}$ 收集逐边不均匀和通道投影误差。对一般实线性场，完整相互作用还包含电子上移同时产生光子、电子下移同时湮灭光子的反旋支路；它们在相互作用绘景中带有约 $2\omega_0$ 的快速相位。这些快速支路能否舍弃，不能只由“$\omega_0$ 很大”判断。前章\eqnrefs{eq:rwa_projected_dynamical_error} 已把旋波余项的真实有限时间误差写成 Duhamel 范数；在脉冲包络缓变且反旋相位频率约为 $2\omega_0$ 的常见情形，分部积分才进一步给出\eqnrefs{eq:rwa-finite-time-error-scales-dp68} 那类充分尺度条件。只有该投影传播误差满足
\[
\epsilon_{\rm RWA}^{\rm dyn}\ll1
\]
时，反旋支路才可受控消去；$|\mathcal G|/\omega_0$ 与 $(\omega_0T)^{-1}$ 只是缓变脉冲下对这个更一般误差界的估计。

保留真实矩阵元的记号 $\mathcal G(t)$，并另定义 $\upsilon(t)=-\ii\mathcal G(t)$，即 $\mathcal G(t)=\ii\upsilon(t)$。这样 $g_{\mathrm Q}=\int\upsilon(t)\dd t$ 与经典 $\beta$ 的相位对应是显式的。代入上式，
\[
\boxed{
\hat H_Q(t)
=
\ii\hbar
\left[
\upsilon(t)\hat b^\dagger\hat a
-
\upsilon^*(t)\hat b\hat a^\dagger
\right].
}
\]
这就是正文\eqnrefs{eq:qpinem-single-mode-hamiltonian-dp68}。
\end{derivation}
